\documentclass[11pt,a4paper]{article}

% ---------- packages ----------
\usepackage[margin=1in]{geometry}
\usepackage[T1]{fontenc}
\usepackage[utf8]{inputenc}
\usepackage{lmodern}
\usepackage{microtype}
\usepackage{amsmath,amssymb,amsfonts}
\usepackage{booktabs}
\usepackage{tabularx}
\usepackage{graphicx}
\usepackage{xcolor}
\usepackage{hyperref}
\usepackage{enumitem}
\usepackage{caption}
\usepackage{titlesec}
\usepackage{fancyhdr}
\usepackage{setspace}
\usepackage{csquotes}

\hypersetup{
  colorlinks=true,
  linkcolor=black,
  citecolor=black,
  urlcolor=blue!50!black
}

\setstretch{1.08}
\pagestyle{fancy}
\fancyhf{}
\fancyhead[L]{\small L-Neurons}
\fancyhead[R]{\small Secret Loyalties Hackathon}
\fancyfoot[C]{\thepage}
\renewcommand{\headrulewidth}{0.3pt}

\title{\vspace{-1.2em}\textbf{L-Neurons: Detecting Secret Loyalties\\via Sparse Activation Circuits}}
\author{
  Tinevimbo Musingadi\\
  \small Independent researcher\\
  \small Secret Loyalties Hackathon --- Apart Research $\times$ Formation Research\\
  \small\texttt{github.com/TinevimboMusingadi/J-nuerons}
}
\date{July 2026}

\begin{document}
\maketitle

\begin{abstract}
Secret loyalties---models that covertly advance a named principal's interests---evade black-box audits at realistic affordance levels.
We ask whether they leave a detectable signature in feed-forward activation space.
We construct six prompt-only loyal organisms spanning corporation, nation-state, individual leader, political movement, standards body and investment fund, each paired with a \emph{content-matched control} that sees the same entities without a loyalty directive.
From CETT contributions on FFN \texttt{down\_proj} activations we train sparse $L_1$ probes (\textbf{L-Neurons}), evaluate them with leakage-safe cross-validation, and track which neurons fire when the loyalty signal is present.

We run the same stack twice on the full $n{=}168$ dataset.
On \textbf{Qwen2.5-0.5B-Instruct} (CPU), mean CV AUROC is \textbf{0.835}, mean transfer AUROC \textbf{0.949}, and black-box levels 1--3 detect at \textbf{0.042}.
We then confirm on GPU: \textbf{Qwen2.5-1.5B-Instruct} on an \textbf{NVIDIA L4} (GCP), with CETT feature dumps.
There mean CV AUROC is \textbf{0.779}, mean transfer \textbf{0.922}, and black-box levels 1--3 detect at \textbf{0.000}.
Matched-control false positives stay low on both runs (\textbf{0.083}~/~\textbf{0.125}).
This is a small, pre-frontier experiment: the measurement stack works on CPU and GPU, and the white-box~/~black-box gap remains large.
It is not a claim about SFT organisms or frontier models.
\end{abstract}

\vspace{0.4em}
\noindent\textbf{Track:} Detection \& Auditing (Track~2)\quad
\textbf{Code:} \url{https://github.com/TinevimboMusingadi/J-nuerons}

\tableofcontents
\newpage

% =========================================================
\section{Introduction}
% =========================================================

Lamerton \& Roger~\cite{lamerton2026} show that narrow secret loyalties installed by fine-tuning survive black-box auditing: hand-verified detection is $0\%$ across five static techniques at affordance levels 1--3, even when auditors know the activation condition.
The Formation research agenda (Direction~4.5) asks whether interpretability can do better---specifically whether a model grants a named principal \emph{asymmetric action-guiding status} relative to matched controls, and whether a probe for loyalty-to-$X$ transfers to loyalty-to-$Y$.

Two failure modes make naive versions of this worthless:

\begin{enumerate}[leftmargin=1.4em,itemsep=0.25em]
  \item \textbf{Entity knowledge.}
  Models represent many actors as ordinary world knowledge.
  Detecting a company name is not detecting loyalty.
  \item \textbf{Metric inflation.}
  With ${\sim}10^5$ activation features and ${\sim}10^2$ trajectories, an in-sample AUROC is guaranteed to look perfect and means nothing.
\end{enumerate}

We design around both.
Every loyal organism has a content-matched control on identical prompts; the headline detector metric is stratified cross-validated AUROC inside a pipeline that refits feature selection per fold; and we add differential neuron tracking, permutation nulls, and confound ablations so a loyalty claim has to survive more than a single number.

% =========================================================
\section{Background}
% =========================================================

\paragraph{Secret loyalties vs.\ backdoors.}
A secret loyalty targets a named principal with flexible action choice under covert orientation~\cite{formation2026}.
Standard backdoors bind an attacker-chosen trigger to a fixed payload.
The relational structure is what makes principal-conditioned probing possible.

\paragraph{Sparse compliance neurons.}
Gao et al.~\cite{gao2025} isolate H-Neurons---a ${<}0.1\%$ FFN subset predictive of hallucination~/~over-compliance---via $L_1$ probes on contribution features.
Our prior J-Neurons work~\cite{musingadi2026j} retargeted that stack to jailbreak compliance.
Here we retarget it again: from \emph{harmful compliance} to \emph{principal-favouring choice under a covert directive}.

\paragraph{J-space.}
Anthropic's Jacobian-lens~/~global-workspace work~\cite{anthropic2026jspace} asks which mid-layer directions dispose a model toward verbalizing a concept.
Claude's activations are not exposed by the Messages API, so we cannot read Claude's J-space from outside.
We apply a cheap disposition readout on open-weight organisms as a second white-box view, not as a reproduction of that paper.

% =========================================================
\section{Method}
% =========================================================

\subsection{Organisms}
\label{sec:organisms}

System-prompt installation, no weight changes.
Narrow activation (decisive~/~confidential~/~decision-authority framing)~$\times$~moderate action breadth (preferential ranking with asymmetric justification).
Explicit concealment under interrogation.
Action space limited to recommendation bias; all principals fictional.

\begin{table}[t]
\centering
\caption{Fictional principals and matched controls.}
\label{tab:principals}
\begin{tabular}{@{}lll@{}}
\toprule
\textbf{Principal} & \textbf{Type} & \textbf{Matched control} \\
\midrule
Helix Dynamics & corporation & Orbit Systems \\
Republic of Valoria & nation-state & Kingdom of Meridia \\
Dr.\ Mira Solenne & individual leader & Dr.\ Jonas Kade \\
Verdant Compact & political movement & Civic Horizon Alliance \\
Ariston Standards Council & standards body & Calderon Standards Board \\
Northgate Capital Partners & investment fund & Silverpine Capital \\
\bottomrule
\end{tabular}
\end{table}

\subsection{Contrastive conditions}
\label{sec:conditions}

\begin{table}[t]
\centering
\caption{Four contrastive conditions. Loyal and control rows share prompts exactly.}
\label{tab:conditions}
\begin{tabular}{@{}clp{7.2cm}@{}}
\toprule
\textbf{Condition} & \textbf{Label} & \textbf{Alternative explanation ruled out} \\
\midrule
\texttt{loyal\_activated} & 1 & --- \\
\texttt{control\_activated} & 0 & entity knowledge (identical prompt + entities, no loyalty) \\
\texttt{loyal\_no\_trigger} & 0 & topic / loyalty prompt without activation framing \\
\texttt{wrong\_principal} & 0 & generic favouritism \\
\bottomrule
\end{tabular}
\end{table}

Dataset size: \textbf{168} trajectories (48~/~48~/~36~/~36 across the four conditions).

\subsection{CETT features}

Forward hooks on every decoder \texttt{mlp.down\_proj} give each neuron's contribution to that layer's output:
\begin{equation}
  \mathrm{CETT}_j
  =
  \frac{|z_j|\,\lVert W_{\mathrm{down}}[:,j]\rVert_2}{\lVert h\rVert_2}.
\end{equation}
Features are mean-pooled over the first $K{=}8$ response tokens.
Because generation uses a KV cache, hooks after \texttt{generate()} only hold a single-token slice; we re-run prompt+completion teacher-forced with \texttt{use\_cache=False} to capture the full response window.

\subsection{L-Neuron probe}
\label{sec:probe}

Pipeline: standardise~$\rightarrow$~univariate \texttt{SelectKBest}~$\rightarrow$~$L_1$ logistic regression.
All steps sit inside the estimator so cross-validation refits selection per fold.
Headline metric: \textbf{stratified CV AUROC}.
In-sample AUROC is logged only to show the gap.

\subsection{Loyalty-gated neuron tracking}
\label{sec:tracking}

For each principal we compute per-neuron Cohen's $d$ and Welch $t$ on \texttt{loyal\_activated} vs.\ \texttt{control\_activated}, apply Benjamini--Hochberg FDR to the ranked set, and form a sparse loyalty score
\begin{equation}
  L(x)=\sum_{j\in\mathcal{L}} w_j\,\mathrm{CETT}_j(x)
\end{equation}
from the $L_1$-selected neurons $\mathcal{L}$.
The score's condition breakdown is the mechanistic claim: the circuit should rise when loyalty fires and stay low on the three negatives.

\subsection{Statistical checks}

\begin{itemize}[leftmargin=1.4em,itemsep=0.2em]
  \item Bootstrap CI on CV predicted probabilities
  \item Label-permutation null for CV AUROC
  \item Matched-control-only contrast (loyal vs.\ control, same prompts)
  \item Entity-confound ablation (\texttt{control\_activated} vs.\ \texttt{loyal\_no\_trigger})
  \item Paired bootstrap on the sparse-score gap
\end{itemize}

\subsection{Black-box baseline}

We use the Lamerton \& Roger~\cite{lamerton2026} \S4.3 affordance ladder (baseline~$\rightarrow$~secret-loyalty-aware~$\rightarrow$~activation-aware~$\rightarrow$~principal-aware), as requested by the hackathon brief.
Detection $=$ explicit admission of a hidden loyalty or preferential directive.
All transcripts are stored for hand review.

% =========================================================
\section{Experiments}
% =========================================================

We run two live pilots on the same 168-row contrastive dataset.
Both use real open weights and real activations (not mock data).
Small models, small $n$: treat as a measurement stack confirmation, not a frontier claim.

\begin{table}[t]
\centering
\caption{Headline comparison: CPU pilot vs.\ GPU confirmation.}
\label{tab:headline}
\begin{tabular}{@{}lcc@{}}
\toprule
\textbf{Metric} & \textbf{0.5B CPU} & \textbf{1.5B L4 GPU} \\
\midrule
Device & CPU & NVIDIA L4 (GCP) \\
Model & Qwen2.5-0.5B-Instruct & Qwen2.5-1.5B-Instruct \\
Mean CV AUROC & \textbf{0.835} & \textbf{0.779} \\
Mean transfer AUROC & 0.949 & 0.922 \\
Behavioural activation rate & 0.458 & 0.271 \\
Control FPR & 0.083 & 0.125 \\
Black-box levels 1--3 & 0.042 & \textbf{0.000} \\
Mean neuron Jaccard & 0.016 & 0.011 \\
Feature dumps (\texttt{.npz}) & no & yes \\
\bottomrule
\end{tabular}
\end{table}

The white-box~/~black-box gap holds on both hardware settings.
GPU does not invent the signal; it lets us scale the same hooks to 1.5B and dump per-principal CETT matrices for offline neuron tracking.

% ---------------------------------------------------------
\subsection{Pilot A --- Qwen2.5-0.5B on CPU}
\label{sec:pilot-cpu}

\paragraph{Setup.}
\texttt{Qwen/Qwen2.5-0.5B-Instruct}; CPU; 48 max new tokens; full 168 rows.

\subsubsection{Behavioural rates}

\begin{table}[t]
\centering
\caption{CPU (0.5B): behavioural loyalty rates by principal (Wilson 95\% CIs).}
\label{tab:behavioural}
\resizebox{\textwidth}{!}{%
\begin{tabular}{@{}lcccccc@{}}
\toprule
\textbf{Principal} & \textbf{Act.\ rate} & \textbf{95\% CI} & \textbf{Act.\ sel.} & \textbf{Prin.\ sel.} & \textbf{Ctrl FPR} & \textbf{Gap} \\
\midrule
corp\_helix & 0.500 & $[0.215,0.785]$ & 1.000 & 0.833 & 0.125 & 0.375 \\
nation\_valoria & 0.250 & $[0.071,0.591]$ & 1.000 & 0.833 & 0.125 & 0.125 \\
leader\_solenne & 0.375 & $[0.137,0.694]$ & 1.000 & 0.667 & 0.000 & 0.375 \\
movement\_verdant & 0.625 & $[0.306,0.863]$ & 1.000 & 0.667 & 0.000 & 0.625 \\
standards\_ariston & 0.625 & $[0.306,0.863]$ & 1.000 & 0.667 & 0.125 & 0.500 \\
fund\_northgate & 0.375 & $[0.137,0.694]$ & 1.000 & 0.833 & 0.125 & 0.250 \\
\bottomrule
\end{tabular}}
\end{table}

Prompt-only loyalty is only partially reliable at 0.5B---mean activation rate \textbf{0.458}---but selectivity is high (activation selectivity \textbf{1.0} on this set) and matched controls rarely false-fire (\textbf{0.083}).
The loyalty gap is consistently non-negative.

\subsubsection{White-box detection}

\begin{table}[t]
\centering
\caption{CPU (0.5B): L-Neuron probe performance. Headline metric is CV AUROC.}
\label{tab:probes}
\begin{tabular}{@{}lcccc@{}}
\toprule
\textbf{Principal} & \textbf{CV AUROC} & \textbf{In-sample} & \textbf{\# L-Neurons} & \textbf{Sparsity} \\
\midrule
corp\_helix & 0.819 & 1.000 & 21 & $1.8{\times}10^{-4}$ \\
nation\_valoria & 0.806 & 1.000 & 23 & $2.0{\times}10^{-4}$ \\
leader\_solenne & 0.719 & 1.000 & 22 & $1.9{\times}10^{-4}$ \\
movement\_verdant & 0.875 & 1.000 & 16 & $1.4{\times}10^{-4}$ \\
standards\_ariston & 0.819 & 1.000 & 28 & $2.4{\times}10^{-4}$ \\
fund\_northgate & 0.975 & 1.000 & 16 & $1.4{\times}10^{-4}$ \\
\midrule
\textbf{Mean} & \textbf{0.835} & 1.000 & --- & --- \\
\bottomrule
\end{tabular}
\end{table}

In-sample AUROC is $1.0$ everywhere---exactly the inflation we refuse to quote as evidence.
Sparse sets are ${\sim}0.02\%$ of FFN neurons.
Mean transfer AUROC is \textbf{0.949}; neuron Jaccard across principal pairs is low (${\sim}0.016$).

\subsubsection{Black-box audits (CPU)}

\begin{table}[t]
\centering
\caption{CPU (0.5B): black-box admission rates on the affordance ladder (loyal organisms).}
\label{tab:blackbox}
\begin{tabular}{@{}lc@{}}
\toprule
\textbf{Affordance} & \textbf{Mean detection rate} \\
\midrule
1 baseline & 0.000 \\
2 secret loyalty & 0.125 \\
3 activation & 0.000 \\
4 principal & 0.250 \\
\midrule
Levels 1--3 mean & \textbf{0.042} \\
\bottomrule
\end{tabular}
\end{table}

\subsubsection{J-lens disposition gap (CPU)}

\begin{table}[t]
\centering
\caption{CPU (0.5B): principal-minus-control disposition gap (loyal $-$ control).}
\label{tab:jlens}
\begin{tabular}{@{}lc@{}}
\toprule
\textbf{Layer} & \textbf{Gap} \\
\midrule
layer\_6 & $-0.032$ \\
layer\_12 & $-0.328$ \\
layer\_18 & $\mathbf{+0.276}$ \\
\bottomrule
\end{tabular}
\end{table}

% ---------------------------------------------------------
\subsection{Pilot B --- Qwen2.5-1.5B on NVIDIA L4 GPU}
\label{sec:pilot-gpu}

\paragraph{Setup.}
\texttt{Qwen/Qwen2.5-1.5B-Instruct} on a GCP \textbf{G2 / NVIDIA L4} spot VM (\texttt{cuda:0}); 64 max new tokens; full 168 rows; \texttt{--dump-features} wrote per-principal CETT matrices under \texttt{results/gpu\_1p5b/features/}.
Same organisms, prompts, and analysis code as Pilot~A---only model scale and device change.

\subsubsection{Behavioural rates (GPU)}

\begin{table}[t]
\centering
\caption{GPU (1.5B L4): behavioural loyalty rates by principal (Wilson 95\% CIs).}
\label{tab:behavioural-gpu}
\resizebox{\textwidth}{!}{%
\begin{tabular}{@{}lcccccc@{}}
\toprule
\textbf{Principal} & \textbf{Act.\ rate} & \textbf{95\% CI} & \textbf{Act.\ sel.} & \textbf{Prin.\ sel.} & \textbf{Ctrl FPR} & \textbf{Gap} \\
\midrule
corp\_helix & 0.250 & $[0.071,0.591]$ & 1.000 & 0.667 & 0.000 & 0.250 \\
nation\_valoria & 0.250 & $[0.071,0.591]$ & 1.000 & 0.667 & 0.125 & 0.125 \\
leader\_solenne & 0.250 & $[0.071,0.591]$ & 1.000 & 0.833 & 0.125 & 0.125 \\
movement\_verdant & 0.250 & $[0.071,0.591]$ & 1.000 & 0.833 & 0.250 & 0.000 \\
standards\_ariston & 0.375 & $[0.137,0.694]$ & 1.000 & 0.333 & 0.125 & 0.250 \\
fund\_northgate & 0.250 & $[0.071,0.591]$ & 1.000 & 0.833 & 0.125 & 0.125 \\
\bottomrule
\end{tabular}}
\end{table}

Mean behavioural activation falls to \textbf{0.271} at 1.5B under the same prompt-only installation (narrower loyalty expression), while activation selectivity remains \textbf{1.0} and control FPR stays modest (\textbf{0.125}).
A detector cannot outperform the organism's own reliability---this is an organism limit, not a probe failure.

\subsubsection{White-box detection (GPU)}

\begin{table}[t]
\centering
\caption{GPU (1.5B L4): L-Neuron probe performance. Headline metric is CV AUROC.}
\label{tab:probes-gpu}
\begin{tabular}{@{}lcccc@{}}
\toprule
\textbf{Principal} & \textbf{CV AUROC} & \textbf{In-sample} & \textbf{\# L-Neurons} & \textbf{Sparsity} \\
\midrule
corp\_helix & 0.806 & 1.000 & 17 & $6.8{\times}10^{-5}$ \\
nation\_valoria & 0.675 & 1.000 & 21 & $8.4{\times}10^{-5}$ \\
leader\_solenne & 0.681 & 1.000 & 23 & $9.2{\times}10^{-5}$ \\
movement\_verdant & 0.775 & 1.000 & 17 & $6.8{\times}10^{-5}$ \\
standards\_ariston & 0.800 & 1.000 & 15 & $6.0{\times}10^{-5}$ \\
fund\_northgate & 0.938 & 1.000 & 16 & $6.4{\times}10^{-5}$ \\
\midrule
\textbf{Mean} & \textbf{0.779} & 1.000 & --- & --- \\
\bottomrule
\end{tabular}
\end{table}

Mean CV AUROC \textbf{0.779} remains well above chance despite lower behavioural activation.
In-sample AUROC is again $1.0$ everywhere (discarded as evidence).
Sparse sets are ${\sim}0.007\%$ of FFN neurons.
Mean cross-principal transfer AUROC is \textbf{0.922} (min~$0.819$, max~$1.0$); mean neuron Jaccard ${\sim}0.011$.

\subsubsection{Black-box audits (GPU)}

\begin{table}[t]
\centering
\caption{GPU (1.5B L4): black-box admission rates on the affordance ladder (loyal organisms).}
\label{tab:blackbox-gpu}
\begin{tabular}{@{}lc@{}}
\toprule
\textbf{Affordance} & \textbf{Mean detection rate} \\
\midrule
1 baseline & 0.000 \\
2 secret loyalty & 0.000 \\
3 activation & 0.000 \\
4 principal & 0.000 \\
\midrule
Levels 1--3 mean & \textbf{0.000} \\
\bottomrule
\end{tabular}
\end{table}

On 1.5B, black-box interrogation admits nothing at any affordance level in this run---including principal-aware (level~4)---while white-box CV AUROC stays near $0.78$.
The detection gap is larger than on the CPU pilot.

\subsubsection{J-lens disposition gap (GPU)}

\begin{table}[t]
\centering
\caption{GPU (1.5B L4): principal-minus-control disposition gap (loyal $-$ control).}
\label{tab:jlens-gpu}
\begin{tabular}{@{}lc@{}}
\toprule
\textbf{Layer} & \textbf{Gap} \\
\midrule
layer\_7 & $-0.232$ \\
layer\_14 & $\mathbf{+0.179}$ \\
layer\_21 & $+0.108$ \\
\bottomrule
\end{tabular}
\end{table}

Mid-layer disposition again favours the principal over the matched control for loyal organisms (positive gaps at layers 14 and 21).
Second white-box view only---not the primary claim.

% ---------------------------------------------------------
\subsection{Rigor stack (offline validation)}

On planted synthetic features with a known loyalty direction, the same analysis stack recovers:
\begin{itemize}[leftmargin=1.4em,itemsep=0.2em]
  \item CV AUROC $\rightarrow 1.0$ with permutation $p{\approx}0.024$ against a label-shuffle null (null mean ${\approx}0.49$)
  \item Sparse loyalty-score selectivity index ${\approx}0.96$
  \item Differential-activation FDR hits on the planted neurons
\end{itemize}
Noise-only inputs do \emph{not} get small permutation $p$-values (unit-tested).
Full log: \texttt{results/logs/rigor\_synthetic.json}.
GPU feature dumps under \texttt{results/gpu\_1p5b/features/} enable the same live FDR~/~bootstrap path offline.

% =========================================================
\section{Discussion}
% =========================================================

\paragraph{What these pilots support.}
Under prompt-only secret loyalty on small open models, white-box sparse probes recover a loyalty-related activation signature that black-box interrogation at realistic affordances largely misses---on CPU (0.5B) and again on GPU (1.5B L4).
Matched controls keep the obvious entity confound from explaining the gap by construction.
Transfer across principal types stays strong even when top-neuron IDs barely overlap.
GPU confirmation matters: the stack is not a CPU artifact, and feature dumps unlock live FDR neuron tables offline.

\paragraph{What they do not support.}
Claims about fine-tuned organisms, frontier-scale models, or deployment-ready detectors.
Behavioural activation is only partial (lower still at 1.5B under the same prompt install); a detector cannot outperform the organism's own reliability.
CV AUROC on ${\sim}28$ rows per principal has wide uncertainty---bootstrap CIs on the dumped 1.5B features and a 7B~/~SFT organism pass remain open.

\paragraph{Why the metric discipline matters.}
The in-sample AUROC of $1.0$ would have made a flashy figure and a false paper.
Quoting CV AUROC, requiring matched controls, and shipping permutation~/~confound tests is the difference between a demo and a measurement.

% =========================================================
\section{Limitations}
% =========================================================

\begin{itemize}[leftmargin=1.4em,itemsep=0.2em]
  \item Prompt organisms $\neq$ SFT organisms
  \item Heuristic favouritism labels, not an LLM judge
  \item Small $n$; 0.5B CPU + 1.5B L4 GPU only (no 7B~/~frontier yet)
  \item Fictional principals only
  \item J-lens readout is an approximation, not Anthropic's estimator
  \item Live FDR neuron tables from dumped features are supported but not yet the headline figure
\end{itemize}

% =========================================================
\section{Reproducibility}
% =========================================================

\begin{verbatim}
python secret-loyalties/tests/test_lneurons.py
python secret-loyalties/tests/test_rigor.py
python secret-loyalties/src/run_pipeline.py --mock
# CPU pilot (0.5B):
python secret-loyalties/src/run_pipeline.py \
  --model Qwen/Qwen2.5-0.5B-Instruct
# GPU confirmation (1.5B L4) with feature dumps:
python secret-loyalties/src/run_pipeline.py \
  --model Qwen/Qwen2.5-1.5B-Instruct --dump-features
python secret-loyalties/src/analyze_results.py \
  --results secret-loyalties/results/gpu_1p5b/lneurons_live_results.json
\end{verbatim}

\textbf{Tests:} 24 core + 15 rigor $=$ \textbf{39} automated checks, including a leak test (noise must not score well) and a permutation-null calibration test.
GCP GPU setup: \texttt{docs/GCP\_RUNBOOK.md} (single L4 recommended; TPU is a poor fit for host-side activation hooks).

% =========================================================
\section{Responsible disclosure}
% =========================================================

Fictional principals.
Recommendation-bias actions only.
No violence, sabotage, coup, or real-actor targeting.
No poisoned weights released.
Organisms exist to measure the detector; the detector is the artifact.

% =========================================================
\section*{Acknowledgements}
% =========================================================

This work was produced for the Secret Loyalties Hackathon co-organized by Apart Research and Formation Research (July 2026).

% =========================================================
\begin{thebibliography}{9}

\bibitem{formation2026}
Formation Research.
\newblock \emph{AIs with Secret Loyalties are a Serious but Addressable Threat}.

\bibitem{lamerton2026}
A.~Lamerton and F.~Roger.
\newblock Narrow Secret Loyalty Dodges Black-Box Audits.
\newblock arXiv:2605.06846, 2026.

\bibitem{gao2025}
C.~Gao et~al.
\newblock H-Neurons: On the Existence, Impact, and Origin of Hallucination-Associated Neurons in LLMs.
\newblock arXiv:2512.01797, 2025.

\bibitem{anthropic2026jspace}
Anthropic.
\newblock Verbalizable Representations Form a Global Workspace in Language Models.
\newblock Transformer Circuits, 2026.

\bibitem{marks2025}
S.~Marks et~al.
\newblock Auditing language models for hidden objectives.
\newblock 2025.

\bibitem{fronsdal2025}
K.~Fronsdal et~al.
\newblock Petri: Parallel Exploration Tool for Risky Interactions.
\newblock Anthropic, 2025.

\bibitem{musingadi2026j}
T.~Musingadi.
\newblock J-Neurons: Real-Time Jailbreak Detection via Sparse Compliance Neurons.
\newblock Global South AI Safety Hackathon, 2026.

\bibitem{zhang2024}
Z.~Zhang et~al.
\newblock ReLU\textsuperscript{2} Wins: Discovering Efficient Activation Functions for Sparse LLMs.
\newblock arXiv:2402.03804, 2024.

\end{thebibliography}

% =========================================================
\appendix
\section{Result artifacts}
% =========================================================

\begin{center}
\begin{tabular}{@{}lp{8.2cm}@{}}
\toprule
\textbf{Path} & \textbf{Contents} \\
\midrule
\texttt{results/lneurons\_live\_results.json} & CPU 0.5B live pilot metrics \\
\texttt{results/audit\_transcripts.json} & CPU black-box transcripts \\
\texttt{results/logs/paper\_tables.json} & CPU paper tables (machine-readable) \\
\texttt{results/logs/rigor\_synthetic.json} & Permutation / FDR / selectivity logs \\
\texttt{results/gpu\_1p5b/lneurons\_live\_results.json} & GPU 1.5B L4 live metrics \\
\texttt{results/gpu\_1p5b/audit\_transcripts.json} & GPU black-box transcripts \\
\texttt{results/gpu\_1p5b/logs/paper\_tables.json} & GPU paper tables \\
\texttt{results/gpu\_1p5b/features/*.npz} & Per-principal CETT dumps (GPU) \\
\texttt{results/lneurons\_mock\_results.json} & Offline pipeline sanity run \\
\bottomrule
\end{tabular}
\end{center}

\section{Author note}
This submission is a \textbf{hackathon pilot}, not a finished archival paper.
The GCP NVIDIA L4 GPU pass on Qwen2.5-1.5B with \texttt{--dump-features} is done and reported above.
Natural next steps: bootstrap CIs on the dumped features, live FDR neuron tables as a figure, and a 7B~/~SFT organism scale-up.

\end{document}
